\documentclass[12pt]{exam}

\usepackage{setspace}
\usepackage{listings}
\usepackage[dvipsnames]{xcolor}
\usepackage{graphicx,subfigure,wrapfig}
\usepackage{multirow}
\usepackage{amsmath}
%\usepackage{multicol}

\usepackage[margin=20mm]{geometry}
\usepackage{xepersian}

\settextfont[
  Path = fonts/,
  Extension = .ttf,
  UprightFont = *,
]{XB Niloofar}

\singlespacing
\parindent 0ex

\newcommand{\class}{سیستم‌های بی‌درنگ}
\newcommand{\term}{فاز یک پروژه}
\newcommand{\prof}{}
\newcommand{\college}{}

\definecolor{CodeBack}{RGB}{248,248,248}
\definecolor{CodeBorder}{RGB}{210,210,210}

\lstset{
basicstyle=\ttfamily\footnotesize,
keywordstyle=\textbf,
identifierstyle=,
stringstyle=\ttfamily,
commentstyle=\color{LimeGreen},
numberstyle=\footnotesize,
showstringspaces=false,
keepspaces=true,
columns=fullflexible,
breaklines=true,
breakatwhitespace=false,
frame=single,
framerule=0.3pt,
rulecolor=\color{CodeBorder},
backgroundcolor=\color{CodeBack},
xleftmargin=6pt,
xrightmargin=6pt,
framesep=5pt}
\begin{document}

% These commands set up the running header on the top of the exam pages
\pagestyle{head}
\firstpageheader{}{}{}
\runningheader{صفحه \thepage\ از \numpages}{}{\class}
\runningheadrule

\begin{tabular}{p{.7\textwidth} l}
\multicolumn{2}{c}{\textbf{به نام خدا}}\\
\multirow{2}{*}{\includegraphics[scale=0.2] {images/logo.png}} & \\ \\
&  \textbf{\class}\\
&  \textbf{\term}\\
&  \textbf{\prof}\\ \\
 \textbf{\college} &  \\
\end{tabular}\\

\rule[1ex]{\textwidth}{.1pt}
\textbf{سبحان آقاسی زاده ۴۰۲۱۰۵۶۶۵ - پارمیس حمصیان ۴۰۲۱۷۰۹۱۳}

\rule[1ex]{\textwidth}{.1pt}
\vspace{0pt}

\begin{questions}
\question

\section*{۱. هدف فاز اول}
هدف پروژه طراحی یک چارچوب مبتنی بر \lr{SDN} برای شبکه‌های محاسبات لبه خودرویی
است. در این شبکه، ماشین‌ها وظیفه‌های محاسباتی تولید می‌کنند. هر وظیفه می‌تواند
روی خود ماشین، روی ماشین ارائه‌دهنده، یا روی سرور \lr{edge} پردازش شود. هدف کلی
این است که تاخیر و گم‌شدن بسته‌ها کمتر شود.

در فاز اول قرار نیست بخش \lr{RL} یا کنترل‌کننده \lr{SDN} کامل را پیاده‌سازی کنیم.
در این فاز فقط پایه شبیه‌سازی ساخته شده است: محیط دوبعدی، پارامترهای دو سناریو،
جایگذاری رندوم ماشین‌های ارائه‌دهنده، و یک روش ساده برای بارسپاری وظیفه. این
فاز نقش روش پایه را دارد تا در فاز دوم بتوانیم روش \lr{SDN/PPO} را با آن مقایسه کنیم.

\section*{۲. سناریوهای شبیه‌سازی}
دو سناریوی خواسته‌شده در صورت پروژه پیاده‌سازی شدند:

\begin{center}
\begin{tabular}{|c|c|c|c|}
\hline
سناریو & تعداد خودروها & تعداد ماشین‌های ارائه‌دهنده & تعداد سرورهای لبه \\
\hline
سناریوی ۱ & ۲۰ & ۳ & ۲ \\
\hline
سناریوی ۲ & ۵۰ & ۶ & ۴ \\
\hline
\end{tabular}
\end{center}

محیط شبیه‌سازی یک صفحه دوبعدی با ابعاد \lr{1000m $\times$ 1000m} است. ماشین‌ها
در این محیط حرکت می‌کنند و وظیفه تولید می‌کنند. برای هر لینک، با توجه به فاصله،
ترافیک و آب‌وهوا، احتمال گم‌شدن بسته محاسبه می‌شود.

\section*{۳. اجزای اصلی پیاده‌سازی}
کد پروژه به چند فایل ساده تقسیم شده است:

\begin{latin}
\begin{lstlisting}
phase1/config.py       scenario and simulation parameters
phase1/models.py       dataclasses for vehicles, tasks, links, metrics
phase1/environment.py  mobility, weather, traffic, packet-loss model
phase1/placement.py    random provider-vehicle placement
phase1/offloading.py   task-offloading and route-selection algorithm
phase1/simulation.py   simulation loop and CSV writers
run_phase1.py          command-line runner
\end{lstlisting}
\end{latin}

این تقسیم‌بندی باعث می‌شود فاز دوم راحت‌تر اضافه شود. مثلا روش رندوم جایگذاری
در فایل \lr{\texttt{placement.py}} است و بعدا می‌توان آن را با کنترل‌کننده
\lr{SDN/PPO} عوض کرد، بدون اینکه کل شبیه‌ساز تغییر کند.

\section*{۴. مدل محیط}
در هر گام زمانی، هر ماشین ممکن است یک وظیفه جدید تولید کند. هر وظیفه سه ویژگی اصلی
دارد:

\begin{itemize}
\item اندازه داده ورودی بر حسب مگابیت
\item حجم پردازش مورد نیاز بر حسب میلیارد سیکل
\item مهلت زمانی انجام وظیفه
\end{itemize}

برای مدل‌سازی شرایط شبکه، چند عامل ساده در نظر گرفته شده است:

\begin{itemize}
\item فاصله بین فرستنده و گیرنده
\item تراکم ترافیک محلی
\item وضعیت آب‌وهوا
\end{itemize}

اگر فاصله دو گره از برد بیشتر باشد، آن لینک نامعتبر است. اگر لینک معتبر باشد،
احتمال گم‌شدن بسته با این رابطه ساده حساب می‌شود:

\begin{latin}
\begin{lstlisting}
loss = base_loss + distance_effect + traffic_effect + weather_penalty
\end{lstlisting}
\end{latin}

این مدل عمدا ساده نگه داشته شده است، ولی همان چیزهایی را که در صورت پروژه آمده
پوشش می‌دهد: فاصله، ترافیک و آب‌وهوا.

\section*{۵. جایگذاری ماشین‌های ارائه‌دهنده}
در فاز اول، جایگذاری ماشین‌های ارائه‌دهنده به صورت رندوم و یکنواخت در محیط
دوبعدی انجام می‌شود:

\begin{latin}
\begin{lstlisting}
x = random value between 0 and area_width
y = random value between 0 and area_height
\end{lstlisting}
\end{latin}

این همان روش پایه است. در فاز دوم، کنترل‌کننده \lr{SDN} می‌تواند با توجه به
ترافیک، آب‌وهوا و تخمین کیفیت لینک، مکان بهتری برای ماشین‌های ارائه‌دهنده انتخاب کند.

\section*{۶. الگوریتم بارسپاری وظیفه}
برای هر وظیفه، الگوریتم همه انتخاب‌های ممکن را بررسی می‌کند:

\begin{itemize}
\item پردازش محلی روی خود ماشین
\item ارسال مستقیم به ماشین ارائه‌دهنده
\item ارسال مستقیم به سرور \lr{edge}
\item ارسال به سرور \lr{edge} از طریق ماشین ارائه‌دهنده
\end{itemize}

گزینه‌هایی که خارج از برد باشند حذف می‌شوند. این همان \lr{action masking} در
نسخه ساده فاز اول است.

برای هر گزینه معتبر، این مقادیر محاسبه می‌شود:

\begin{itemize}
\item تاخیر ارسال
\item تاخیر پردازش
\item احتمال گم‌شدن بسته
\item تعداد مورد انتظار \lr{retransmission}
\item اینکه \lr{deadline} از دست می‌رود یا نه
\end{itemize}

سپس برای هر گزینه یک امتیاز محاسبه می‌شود:

\begin{latin}
\begin{lstlisting}
score = delay + loss_weight * packet_loss + deadline_penalty
\end{lstlisting}
\end{latin}

گزینه‌ای که کمترین \lr{score} را داشته باشد انتخاب می‌شود. این روش \lr{RL} نیست،
اما برای فاز اول ساده، قابل توضیح و کامل است.

\section*{۷. معیارهای خروجی}
معیارهای زیر ذخیره می‌شوند:

\begin{itemize}
\item میانگین نسبت گم‌شدن بسته
\item میانگین تاخیر سیستم
\item تعداد وظیفه‌هایی که \lr{deadline} را از دست داده‌اند
\item تعداد رخدادهای نقطه کور
\item میانگین موقعیت ماشین‌های ارائه‌دهنده
\item تعداد انتخاب هر مقصد و هر مسیر
\end{itemize}

خروجی‌ها در پوشه \lr{\texttt{outputs/phase1}} و به صورت فایل \lr{CSV} ذخیره می‌شوند.
فایل \lr{\texttt{summary.csv}} خلاصه دو سناریو را دارد و فایل‌های
\lr{\texttt{steps\_scenario\_1.csv}} و \lr{\texttt{steps\_scenario\_2.csv}} معیارهای
هر گام زمانی را ذخیره می‌کنند.

\section*{۸. نتیجه اجرای فعلی}
پس از اجرای پیش‌فرض برنامه با دستور زیر:

\begin{latin}
\begin{lstlisting}
python3 run_phase1.py
\end{lstlisting}
\end{latin}

خلاصه نتیجه به شکل زیر است:

\begin{center}
\begin{tabular}{|c|c|c|c|c|c|}
\hline
سناریو & وظیفه‌ها & میانگین گم‌شدن بسته & میانگین تاخیر & \lr{deadline} از دست‌رفته & نقطه کور \\
\hline
۱ & ۸۲۳ & ۰.۰۹۶۲ & ۰.۴۱۱۸ & ۲۷ & ۱۵۹ \\
\hline
۲ & ۲۰۷۶ & ۰.۱۰۰۸ & ۰.۳۷۴۵ & ۵۳ & ۱۷۶ \\
\hline
\end{tabular}
\end{center}

این اعداد خروجی روش پایه هستند. در فاز دوم می‌توان روش \lr{RL/SDN} را با همین
سناریوها و همین معیارها با این \lr{baseline} مقایسه کرد.

\section*{۹. جمع‌بندی}
در فاز اول، محیط شبیه‌سازی، پارامترهای دو سناریو، جایگذاری رندوم ماشین‌های
ارائه‌دهنده و الگوریتم بارسپاری وظیفه پیاده‌سازی شد. کد وابستگی خارجی ندارد تا
اجرای آن ساده باشد. ساختار پروژه هم طوری نوشته شده که در فاز دوم بتوان
کنترل‌کننده \lr{SDN}، فیلتر کالمن، \lr{RL} و نمودارهای نهایی را روی همین پایه اضافه کرد.

\end{questions}
\end{document}
